\documentclass{article}

% if you need to pass options to natbib, use, e.g.:
%     \PassOptionsToPackage{numbers, compress}{natbib}
% before loading neurips_2021

% ready for submission
%\usepackage{neurips_2021}

% to compile a preprint version, e.g., for submission to arXiv, add add the
% [preprint] option:
%     \usepackage[preprint]{neurips_2021}

% to compile a camera-ready version, add the [final] option, e.g.:
%     \usepackage[final]{neurips_2021}

% to avoid loading the natbib package, add option nonatbib:
%    \usepackage[nonatbib]{neurips_2021}

\usepackage[utf8]{inputenc} % allow utf-8 input
%\usepackage[T1]{fontenc}    % use 8-bit T1 fonts
\usepackage{hyperref}       % hyperlinks
\usepackage{url}            % simple URL typesetting
\usepackage{booktabs}       % professional-quality tables
\usepackage{amsfonts}       % blackboard math symbols
\usepackage{nicefrac}       % compact symbols for 1/2, etc.
\usepackage{microtype}      % microtypography
\usepackage{xcolor}         % colors
\usepackage{fullpage}
\usepackage{authblk}
\usepackage{hyperref,url}
\usepackage{sidecap}
\sidecaptionvpos{figure}{c}

\title{Numerical Composition of Differential Privacy\footnote{Author ordering is alphabetical. Code is available at \url{https://github.com/microsoft/prv_accountant}.}}

% The \author macro works with any number of authors. There are two commands
% used to separate the names and addresses of multiple authors: \And and \AND.
%
% Using \And between authors leaves it to LaTeX to determine where to break the
% lines. Using \AND forces a line break at that point. So, if LaTeX puts 3 of 4
% authors names on the first line, and the last on the second line, try using
% \AND instead of \And before the third author name.





\author[1]{Sivakanth Gopi}
\author[2]{Yin Tat Lee}
\author[1]{Lukas Wutschitz}
\affil[1]{Microsoft\\
\texttt{\{sigopi,lukas.wutschitz\}@microsoft.com}}
\affil[2]{University of Washington\\
\texttt{yintat@uw.edu}}
% \affil[3]{Microsoft\\
% \texttt{lukas.wutschitz@microsoft.com}}

\date{}


%%%%%%%%%%%%%%%%%%%%% MACROS %%%%%%%%%%%%%%%%%%%%%%%%%%%%%%%
\usepackage{amsthm,thmtools}
\usepackage{thm-restate}
\usepackage{amsmath,amssymb}
\usepackage{graphicx,subcaption}
\usepackage{float}
\usepackage{wrapfig}
\usepackage[ruled,vlined]{algorithm2e}

\usepackage{placeins}
\input{stdinc}

\declaretheorem[within=section]{theorem}
\declaretheorem[sibling=theorem]{corollary}
\declaretheorem[sibling=theorem]{assumption}

\declaretheorem[sibling=theorem]{lemma}
\declaretheorem[sibling=theorem]{claim}

\declaretheorem[sibling=theorem]{definition}

\declaretheorem[sibling=theorem]{proposition}
\declaretheorem[sibling=theorem]{remark}
\declaretheorem[sibling=theorem]{question}
\declaretheorem[sibling=theorem]{conjecture}
\declaretheorem[sibling=theorem]{example}
\declaretheorem[sibling=theorem]{fact}

\newcommand{\Lap}{\mathsf{Lap}}
\newcommand{\Gauss}{\mathsf{Gauss}}
\DeclareMathOperator{\E}{\mathbb{E}}

\newcommand{\grad}{\nabla}
%\newcommand{\tg}{\widetilde{g}}
\newcommand{\Id}{\mathrm{Id}}
\newcommand{\SF}{\mathrm{SF}}
\newcommand{\CDF}{\mathrm{CDF}}
\newcommand{\sample}{\mathsf{Sample}}
\newcommand{\ones}{\mathbf{1}}
\newcommand{\tP}{\widetilde{P}}
\newcommand{\tQ}{\widetilde{Q}}
\newcommand{\tZ}{\widetilde{Z}}
\newcommand{\tz}{\widetilde{z}}
\newcommand{\ty}{\widetilde{y}}
\newcommand{\tY}{\widetilde{Y}}
\newcommand{\tx}{\widetilde{x}}
\newcommand{\tX}{\widetilde{X}}


\newcommand{\vjp}{\texttt{vjp}}
\newcommand{\jvp}{\texttt{jvp}}
\graphicspath{{Figures/}}



\newcommand*{\affaddr}[1]{#1} 
\newcommand*{\affmark}[1][*]{\textsuperscript{#1}}
\newcommand*{\email}[1]{\texttt{#1}}


\newcommand{\error}{\mathrm{error}}
\newcommand{\upper}{\mathrm{upper}}
\newcommand{\lp}{\left(}
\newcommand{\lb}{\left[}
\newcommand{\rb}{\right]}
\newcommand{\rp}{\right)}
\newcommand{\bR}{\overline{\R}}

\newcommand{\num}[1]{#1}
\newcommand{\SI}[1]{#1}
\newcommand{\second}{\ sec}

%%%%%%%%%%%%%%%%%%%%%%%%%%%%%%%%%%%%%%%%%%%%%%%%%%%%%%%%%%%%%

%%% disable note. %%%
\iffalse
\let\gnote\undefined
\let\ynote\undefined
\newcommand{\gnote}[1]{}%
\newcommand{\ynote}[1]{}%
\fi
%%%%%%%%%%%%%%%


\begin{document}

\maketitle

\begin{abstract}
We give a fast algorithm to optimally compose privacy guarantees of differentially private (DP) algorithms to arbitrary accuracy. Our method is based on the notion of \emph{privacy loss random variables} to quantify the privacy loss of DP algorithms. %This has an intuitive meaning: if a DP algorithm $M$ has privacy random variables $(X,Y)$ (which are real-valued), then inferring membership of any user from the output of $M$ is at least as hard as distinguishing between $X$ and $Y$ from a random sample. Though mathematically equivalent to the $(\epsilon,\delta)$-DP framework of Dwork et al.~\cite{DMNS06} and the recently proposed $f$-DP framework of Dong et al.~\cite{DongRS19}, this notion allows for a much simpler way to compute the privacy loss of \emph{composition} of DP algorithms, compared to the other frameworks. If $M_1$ and $M_2$ are DP algorithms with privacy random variables $(X_1,Y_1)$ and $(X_2,Y_2)$ respectively, then the (adaptive) composition $M_1\circ M_2$ has privacy random variables $(X_1+X_2,Y_1+Y_2)$.
The running time and memory needed for our algorithm to approximate the privacy curve of a DP algorithm composed with itself $k$ times is $\tilde{O}(\sqrt{k})$. This improves over the best prior method by Koskela et al.~\cite{koskela2021computing} which requires $\tilde{\Omega}(k^{1.5})$ running time. We demonstrate the utility of our algorithm by accurately computing the privacy loss of DP-SGD algorithm of Abadi et al.~\cite{Abadi16} and showing that our algorithm speeds up the privacy computations by a few orders of magnitude compared to prior work, while maintaining similar accuracy.
\end{abstract}

% We give a fast algorithm to optimally compose privacy guarantees of differentially private (DP) algorithms to arbitrary accuracy. Our method is based on the notion of privacy loss random variables to quantify the privacy loss of DP algorithms. The running time and memory needed for our algorithm to approximate the privacy curve of a DP algorithm composed with itself $k$ times is $\tilde{O}(\sqrt{k})$. This improves over the best prior method by Koskela et al. (2020) which requires $\tilde{\Omega}(k^{1.5})$ running time. We demonstrate the utility of our algorithm by accurately computing the privacy loss of DP-SGD algorithm of Abadi et al. (2016) and showing that our algorithm speeds up the privacy computations by a few orders of magnitude compared to prior work, while maintaining similar accuracy.
\newpage
\tableofcontents
\newpage




%\ynote{Need to handle the question 3 and 4 in checklist}
\section{Introduction}
\label{sec:intro}
\input{introduction}

\section{DP Preliminaries}
\label{sec:prelims}
%!TEX root=./MAIN.tex
%The $f$-DP framework, recently proposed in~\cite{DongRS19}, allows us to reason about a collection of $(\eps,\delta)$-privacy guarantees simultaneously which can then be composed to get a much better $(\eps,\delta)$-privacy for the final algorithm. 
%We will first define the notion of $(\eps,\delta)$-DP formally and then define the notion of $f$-DP. We then state a proposition from~\cite{DongRS19} which shows that these two notions are dual to each other.

%We will begin with definition of $(\eps,\delta)$-differentially private algorithm.
%\begin{definition}[$(\eps,\delta)$-DP~\cite{DMNS06}]
%	We say that an algorithm $\cM$ is $(\eps,\delta)$-DP if for any two neighboring databases $D,D'$ and any subset $S$ of outputs, we have $\Pr[\cM(D)\in S] \le e^\eps\Pr[\cM(D')\in S]+\delta.$
%\end{definition}

Given a DP algorithm $\cM$, for each value of $\eps\ge 0$, there exists some $\delta \in [0,1]$ such that $\cM$ is $(\eps,\delta)$-DP. We can represent all these privacy guarantees by a function $\delta_\cM(\eps):\R^{\ge 0} \to [0,1]$ and say that $\delta_\cM(\cdot)$ is the \emph{privacy curve} of $\cM.$ This inspires the following definition of a privacy curve between two random variables.

\begin{definition}[Privacy curve]
	Given two random variables $X,Y$ supported on some set $\Omega$, define $\delta(X||Y): \R \to [0,1]$ as:
	$$\delta(X||Y)(\eps) = \sup_{S\subset \Omega} \Pr[Y\in S] - e^\eps \Pr[X\in S].$$
\end{definition}

Therefore an algorithm $\cM$ is $(\eps,\delta)$-DP iff $\delta\lp \cM(D)||\cM(D')\rp(\eps) \le \delta$ for all neighboring databases $D,D'.$ 

\begin{remark}
	Note that not all functions $\delta:\R \to [0,1]$ are privacy curves. A characterization of privacy curves can be obtained using the $f$-DP framework of~\cite{DongRS19}. The notion of privacy curve $\delta(X||Y)$ and tradeoff function $T(X||Y)$ are dual to each other via convex duality~\cite{DongRS19}. This implies a characterization of privacy curves as shown in \cite{zhu2021optimal}.
\end{remark}

\begin{definition}[Composition of privacy curves~\cite{DongRS19}]
	Let $\delta_1\equiv \delta(X_1||Y_1)$ and $\delta_2\equiv \delta(X_2||Y_2)$ be any two privacy curves. The composition of the privacy curves, denoted by $\delta_1 \otimes \delta_2$, is defined as $$\delta_1 \otimes \delta_2 \equiv \delta\lp(X_1,X_2)||(Y_1,Y_2)\rp$$ where $X_1,X_2$ are independently sampled and $Y_1,Y_2$ are independently sampled.
\end{definition}
Note that there can be many pairs of random variables which have the same privacy curve, but the above operation is well-defined. If $\delta(X_1||Y_1)\equiv \delta(X_1'||Y_1')$ and $\delta(X_2||Y_2)\equiv \delta(X_2'||Y_2')$, then it was shown by~\cite{DongRS19} that $$\delta\lp(X_1,X_2)||(Y_1,Y_2)\rp=\delta\lp(X_1',X_2')||(Y_1',Y_2')\rp.$$ \cite{DongRS19} also show that $\otimes$ is a commutative and associative operation.

Given two DP algorithms $M_1$ and $M_2$, the adaptive composition $(M_2 \circ M_1)(D)$ is an algorithm which outputs $(M_1(D),M_2(D,M_1(D))$, i.e., $M_2$ can look at the database $D$ and also the output of the previous algorithm $M_1(D)$. Adaptive composition of more than two algorithms is similarly defined. Suppose $M_1$ has privacy curve $\delta_1$ and $M_2$ has privacy curve $\delta_2$ (i.e., $M_2(\cdot, y)$ is a DP algorithm with privacy curve $\delta_2$ for any fixed $y.$). The following composition theorem shows how to get the privacy curve of $M_2 \circ M_1.$
\begin{theorem}[Composition theorem~\cite{DongRS19}]
\label{prop:composition}
	Let $M_1,M_2,\dots,M_k$ be DP algorithms with privacy curves given by $\delta_1,\delta_2,\dots,\delta_k$ respectively. The privacy curve of the adaptive composition $M_k\circ M_{k-1} \circ \dots \circ M_1$ is given by $\delta_1\otimes \delta_2 \otimes \cdots \otimes \delta_k.$
\end{theorem}








%\newpage
\section{Privacy Loss Random Variables (PRVs)}
\label{sec:PRV}
%!TEX root=./MAIN.tex

%Given a privacy curve $\delta$, there can be many pairs of random variables $(X,Y)$ such that $\delta(X||Y)\equiv \delta$. 
The notion of \emph{privacy loss random variables} (PRVs) is a unique way to assign a pair $(X,Y)$ for any privacy curve $\delta$ such that $\delta \equiv \delta(X||Y).$ PRVs allow us to compute composition of two algorithms via summing random variables (Theorem~\ref{thm:composition_PRVs}) (equivalently, convolving their distributions). Thus PRVs can be thought of as a \emph{reparametrization of privacy curves} where composition becomes convolution. In this paper, we differ from the usual definition of PRVs given in~\cite{DR16,koskela2020computing}, which are tied to a specific algorithm. Instead we think of them as a reparametrization of privacy curves and study them directly. This allows us to succinctly prove many useful properties of PRVs.

Let $\bR=\R\cup\set{-\infty,\infty}$ be the extended real line where we define $\infty+x=\infty$ and $-\infty+x=-\infty$ for $x\in \R.$


\begin{definition}[Privacy loss random variables (PRVs)]
	Given a privacy curve $\delta:\R\to[0,1]$, we say that $(X,Y)$ are \emph{privacy loss random variables} for $
	\delta$, if they satisfy the following conditions:
	\begin{itemize}
	 	\item $X,Y$ are supported on $\bR$,
	 	\item $\delta(X||Y)\equiv \delta$,
	 	\item $Y(t)=e^t X(t)$ for every $t\in \R$ and
	 	\item $Y(-\infty)=0$ and $X(\infty)=0$
	 \end{itemize} 
	 where $X(t),Y(t)$ are probability density functions of $X,Y$ respectively.
\end{definition}
Mathematically, the correct way to write the condition $Y(t)=e^t X(t)$ is to say that $\E_{Y}[\phi(Y)] = \E_{X}[\phi(X)e^X]$ for all test functions $\phi:\bR \to [0,1]$ with $\phi(\infty)=\phi(-\infty)=0.$ This will generalize to all situations where $X,Y$ are continuous or discrete or both. 
%Equivalently, one can rewrite the condition as $\frac{dY}{dX}(t)=e^t$ where $\frac{dY}{dX}$ is the Radon-Nikodym derivative of the distribution of $Y$ w.r.t to the distribution of $X$. 
For ease of exposition, we ignore this complication and assume that $X(t),Y(t)$ represent the PDFs if $X,Y$ are continuous at $t$, or the probability masses if they have point masses at $t$.




The following theorem shows that the PRVs for a privacy curve $\delta=\delta(P||Q)$ are given by the log-likelihood random variables of $P,Q.$ 
\begin{restatable}{theorem}{PRVLLRV}
\label{thm:PRV_LLRV}
 Let $\delta:\R \to [0,1]$ be a privacy curve given by $\delta \equiv \delta(P||Q)$ where $P,Q$ are two random variables supported on $\Omega$. The PRVs $(X,Y)$ for the privacy curve $\delta$ are given by\footnote{Here $Q(\omega)$ and $P(\omega)$ are the probability density functions of $Q,P$ respectively. Note that the mathematically precise way is to replace the ratio $\frac{Q(\omega)}{P(\omega)}$ by the Radon-Nikodym derivative $\frac{dQ}{dP}(\omega).$}:
 $$X = \log\lp \frac{Q(\omega)}{P(\omega)}\rp \text{ where } \omega \sim P,$$
 $$Y = \log\lp \frac{Q(\omega)}{P(\omega)}\rp \text{ where } \omega \sim Q.$$
 % Moreover, the privacy curve $\delta$ can be expressed in terms of PRVs $(X,Y)$ as:
 % \begin{equation}
 % \label{eqn:delta_PRV_simple}
 % \delta(\eps)=\Pr[Y> \eps] - e^\eps \Pr[X > \eps]=\E_Y[(1-e^{\eps-Y})_+] =\Pr[Y\ge \eps+Z].
 % \end{equation}
 % where $Z$ is an exponential random variable.\footnote{For $x\in \R, x_+=\max\{x,0\}$.}
\end{restatable}

The following theorem provides a formula for computing the privacy curve $\delta$ in terms of the PRVs and conversely a formula for PRVs in terms of the privacy curve. A similar statement appears in~\cite{sommer2019privacy,koskela2020computing}.

\begin{restatable}{theorem}{PRVprivacycurve}
\label{thm:PRV_privacycurve}
 The privacy curve $\delta$ can be expressed in terms of PRVs $(X,Y)$ as:
 \begin{equation}
 \label{eqn:delta_PRV_simple}
 \delta(\eps)=\Pr[Y> \eps] - e^\eps \Pr[X > \eps]=\E_Y[(1-e^{\eps-Y})_+] =\Pr[Y\ge \eps+Z].
 \end{equation}
 where $Z$ is an exponential random variable.\footnote{For $x\in \R, x_+=\max\{x,0\}$.} Conversely, given a privacy curve $\delta:\R\to [0,1]$, we can compute the PDFs of its PRVs $(X,Y)$ as:
 \begin{equation}
 \label{eqn:PRV_from_delta}
 Y(t) = \delta''(t)-\delta'(t) \text{ and } X(t)=e^t(\delta''(t)-\delta'(t)).
 \end{equation}
\end{restatable}

\begin{remark}
	Theorem~\ref{thm:PRV_privacycurve} shows that the PRVs $X,Y$ do not depend on the particular $P,Q$ used to represent the privacy curve $\delta$ in Theorem~\ref{thm:PRV_LLRV}. So we should think of the PDF of of the PRV $Y$ (or $X$) as an equivalent reparametrization of the privacy curve $\delta:\R \to [0,1]$, just as the notion of $f$-DP~\cite{DongRS19} is a reparametrization of the privacy curve $\delta$. %In particular any random variable $Y$ supported on $\R\cup {\infty}$ which satisfies $\E[e^{-Y}]=1$ is the PRV of some privacy curve.\footnote{This condition arises because $X(t)=e^tY(t)$ and $\E[e^{-Y}]=\int_{-\infty}^{\infty} e^{-t}Y(t) dt =\int_{-\infty}^{\infty} X(t) dt =1$.}
\end{remark}





PRVs are useful in computing privacy curves because the composition of two privacy curves can be computed by adding the corresponding pairs of PRVs. A similar statement appears in~\cite{DR16}.
\begin{restatable}{theorem}{compositePRV}
\label{thm:composition_PRVs}
Let $\delta_1,\delta_2$ be two privacy curves with PRVs $(X_1,Y_1)$ and $(X_2,Y_2)$ respectively. Then the PRVs for $\delta_1 \otimes \delta_2=\delta(X_1,X_2||Y_1,Y_2)$ are given by $(X_1+X_2,Y_1+Y_2)$. In particular, $$\delta_1\otimes \delta_2 = \delta(X_1+X_2||Y_1+Y_2).$$
\end{restatable}
\begin{proof}
Let $(X,Y)$ be the privacy random variables for $\delta(X_1,X_2||Y_1,Y_2)$. By Theorem~\ref{thm:PRV_LLRV},
\begin{align*}
	X & = \log\lp\frac{(Y_1,Y_2)(t_1,t_2)}{(X_1,X_2)(t_1,t_2)}\rp \text{ where } (t_1,t_2)\sim (X_1,X_2)\\
	& = \log\lp\frac{Y_1(t_1)Y_2(t_2)}{X_1(t_1)X_2(t_2)}\rp \text{ where } t_1\sim X_1,t_2\sim X_2\tag{By independence of $X_1,X_2$ and indpendence of $Y_1,Y_2$}\\
	& = \log\lp e^{t_1} \cdot e^{t_2}\rp \text{ where } t_1\sim X_1,t_2\sim X_2\\
	& = t_1+t_2 \text{ where } t_1\sim X_1,t_2\sim X_2\\
	& = X_1 + X_2.
\end{align*}
Similarly, \begin{align*}
	Y & = \log\lp\frac{(Y_1,Y_2)(t_1,t_2)}{(X_1,X_2)(t_1,t_2)}\rp \text{ where } (t_1,t_2)\sim (Y_1,Y_2)\\
	& = t_1+t_2 \text{ where } t_1\sim Y_1,t_2\sim Y_2\\
	& = Y_1 + Y_2.
\end{align*}
\end{proof}


In Appendix~\ref{sec:appendix_PRV}, we provide a proof of Theorems \ref{thm:PRV_LLRV} and \ref{thm:PRV_privacycurve}. We also discuss how to compute the PRVs for a subsampled mechanism given the PRVs for the original mechanism and give examples of PRVs for few standard mechanisms. These are used in our experiments to calculate the PRVs for DP-SGD.



\section{Numerical composition of privacy curves}
\label{sec:numerical_composition}
\input{numerical_composition.tex}

\section{Experiments}
\label{sec:experiments}
\input{experiments.tex}


\subsection*{Acknowledgements}
We would like to thank Janardhan Kulkarni and Sergey Yekhanin for several useful discussions and encouraging us to work on this problem.
L.W. would like to thank Daniel Jones and Victor R\"uhle for fruitful discussions and helpful guidance.

\FloatBarrier
\bibliographystyle{alpha}
\bibliography{references}




\appendix
\input{appendix_num_acc}
%!TEX root=./MAIN.tex
\section{Privacy Loss Random Variables}
\label{sec:appendix_PRV}
In this section, we continue the discussion on privacy random variables in Section \ref{sec:PRV}. First, we give the proof of the formula for PRVs of $\delta(P||Q)$ and the formula for a privacy curve given its PRVs (Theorem \ref{thm:PRV_LLRV}).

{\renewcommand\footnote[1]{}\PRVLLRV*}
\begin{proof}
We will first verify that $Y(t)=e^t X(t).$ This is equivalent to proving that $\E_Y[\phi(Y)]=\E_X[\phi(X)e^X]$ for any test function $\phi:\bR \to [0,1]$. This is true since
	\begin{align*}
		\E_Y[\phi(Y)]&= \E_{\omega\sim Q}\lb \phi\lp\log\lp \frac{Q(\omega)}{P(\omega)}\rp\rp\rb \\
		&= \E_{\omega\sim P}\lb \phi\lp\log\lp \frac{Q(\omega)}{P(\omega)}\rp\rp \frac{Q(\omega)}{P(\omega)}\rb\\
		&= \E_X[\phi(X)e^X].
	\end{align*}
	We will now prove that $\delta(X||Y)=\delta(P||Q).$
	We have
	\begin{align*}
		\delta(P||Q)(\eps) &= \sup_{S \subset \Omega} \Pr[Q\in S] - e^\eps\Pr[P\in S]\\
		&= \Pr[Q\in S_\eps] - e^\eps\Pr[P\in S_\eps]
	\end{align*}
	where $$S_\eps = \left\{\omega\in \Omega: \frac{Q(\omega)}{P(\omega)} > e^\eps \right\}=\left\{\omega\in \Omega: \log\lp\frac{Q(\omega)}{P(\omega)}\rp > \eps \right\}.$$ 
	Therefore $\Pr[Q\in S_\eps]=\Pr[Y > \eps]$ and $\Pr[P\in S_\eps]=\Pr[X > \eps]$. To complete the proof, note that 
	\begin{align*}
		\delta(X||Y)(\eps) &= \sup_{T \subset \bR} \Pr[Y\in T] - e^\eps\Pr[X\in T]\\
		&= \Pr[Y\in T_\eps] - e^\eps\Pr[X\in T_\eps]
	\end{align*}
	where $$T_\eps = \left\{t\in \bR: \frac{Y(t)}{X(t)} > e^\eps \right\}= \left\{t\in \bR: e^t > e^\eps \right\}=(\eps,\infty].$$
	Putting it all together, we have:
	\begin{align*}
		\delta(P||Q)(\eps)&=\Pr[Y>\eps] - e^\eps \Pr[X>\eps]=\delta(X||Y).
	\end{align*}
% Finally, since the PDFs of PRVs $(X,Y)$ satisfy the relation $Y(t)=e^tX(t)$, we can rewrite the equation~\ref{eqn:delta_PRV_simple} in terms of just $Y$ or just $X$.
% 		\begin{align*}
% 		\delta(\eps)&=\Pr[Y\ge \eps]-e^\eps \Pr[X\ge \eps]\\
% 		&=\int_\eps^\infty Y(t) dt - \int_\eps^\infty e^\eps X(t) dt\\
% 		&=\int_\eps^\infty Y(t) dt - \int_\eps^\infty e^{\eps-t} Y(t) dt \tag{Since $Y(t)=e^tX(t)$}\\
% 		&=\int_\eps^\infty Y(t) (1-e^{\eps-t}) dt\\
% 		&=\int_{-\infty}^\infty Y(t) (1-e^{\eps-t})_+ dt\\
% 		&=\E_Y[(1-e^{\eps-Y})_+]
% 	\end{align*}
% 	To get the other form for $\delta(\eps)$, we use the integration by parts formula.
% 	\begin{align*}
% 		\delta(\eps)&=\int_{\eps}^\infty Y(t) (1-e^{\eps-t}) dt\\
% 		&=\int_{\eps}^\infty Y(t) dt + \int_{\eps}^\infty \lp -Y(t) \rp e^{\eps-t} dt\\
% 		&=\Pr[Y\ge \eps] + \lp \Pr[Y\ge t]e^{\eps-t}\Big|_\eps^\infty - \int_{\eps}^\infty \Pr[Y\ge t] \lp -e^{\eps-t}\rp dt \rp\\
% 		&=\Pr[Y\ge \eps] - \Pr[Y\ge \eps] + \int_{\eps}^\infty \Pr[Y\ge t] e^{\eps-t} dt\\
% 		&=\int_\eps^\infty e^{\eps-t} \Pr[Y\ge t] dt \\
% 		&=\int_0^\infty e^{-z} \Pr[Y\ge \eps+z] dz  \tag{Substituting $z=t-\eps$}\\
% 		&=\Pr[Y\ge \eps+Z]. \tag{where $Z$ is an exponential random variable}
% 	\end{align*}
\end{proof}

{\renewcommand\footnote[1]{}\PRVprivacycurve*}
\begin{proof}	
Since the PDFs of PRVs $(X,Y)$ satisfy the relation $Y(t)=e^tX(t)$, we can rewrite the equation~\ref{eqn:delta_PRV_simple} in terms of just $Y$ or just $X$.
		\begin{align*}
		\delta(\eps)&=\Pr[Y\ge \eps]-e^\eps \Pr[X\ge \eps]\\
		&=\int_\eps^\infty Y(t) dt - \int_\eps^\infty e^\eps X(t) dt\\
		&=\int_\eps^\infty Y(t) dt - \int_\eps^\infty e^{\eps-t} Y(t) dt \tag{Since $Y(t)=e^tX(t)$}\\
		&=\int_\eps^\infty Y(t) (1-e^{\eps-t}) dt\\
		&=\int_{-\infty}^\infty Y(t) (1-e^{\eps-t})_+ dt\\
		&=\E_Y[(1-e^{\eps-Y})_+]
	\end{align*}
	To get the other form for $\delta(\eps)$, we use the integration by parts formula.
	\begin{align*}
		\delta(\eps)&=\int_{\eps}^\infty Y(t) (1-e^{\eps-t}) dt\\
		&=\int_{\eps}^\infty Y(t) dt + \int_{\eps}^\infty \lp -Y(t) \rp e^{\eps-t} dt\\
		&=\Pr[Y\ge \eps] + \lp \Pr[Y\ge t]e^{\eps-t}\Big|_\eps^\infty - \int_{\eps}^\infty \Pr[Y\ge t] \lp -e^{\eps-t}\rp dt \rp\\
		&=\Pr[Y\ge \eps] - \Pr[Y\ge \eps] + \int_{\eps}^\infty \Pr[Y\ge t] e^{\eps-t} dt\\
		&=\int_\eps^\infty e^{\eps-t} \Pr[Y\ge t] dt \\
		&=\int_0^\infty e^{-z} \Pr[Y\ge \eps+z] dz  \tag{Substituting $z=t-\eps$}\\
		&=\Pr[Y\ge \eps+Z]. \tag{where $Z$ is an exponential random variable}
	\end{align*}
We now prove the converse relation by differentiating the expression for $\delta(\eps)$ twice. We have:
	\begin{align*}
 	&\delta(\eps)=\int_\eps^\infty Y(t) dt - e^{\eps}\int_\eps^\infty e^{-t} Y(t) dt\\
	\implies & \delta'(\eps)=-Y(\eps)+e^{\eps}\cdot e^{-\eps}Y(\eps)-e^{\eps}\cdot \int_\eps^\infty e^{-t} Y(t) dt = -e^{\eps}\cdot \int_\eps^\infty e^{-t} Y(t) dt \\
	\implies & e^{-\eps}\delta'(\eps)=-\int_\eps^\infty e^{-t} Y(t) dt \\
	\implies & e^{-\eps}\delta''(\eps)-e^{-\eps}\delta'(\eps)=e^{-\eps} Y(\eps) \\
	\implies & Y(\eps)=\delta''(\eps)-\delta'(\eps).
	\end{align*}


\end{proof}




\subsection{Examples of privacy loss random variables}
In this section, we state the PRVs for a few standard mechanisms. 

\begin{proposition}[Gaussian Mechanism]
The PRVs for $\delta(\cN(\mu,1)||\cN(0,1))$ are:
$$X=\cN(-\mu^2/2,\mu^2) \text{ and } Y=\cN(\mu^2/2,\mu^2).$$
\end{proposition}
\begin{proof}
	Let $P=\cN(\mu,1)$ and $Q=\cN(0,1)$.
	By Theorem~\ref{thm:PRV_LLRV},
	\begin{align*}
		Y &\sim \log\lp \frac{Q(t)}{P(t)}\rp \text{ where } t\sim Q\\
		&\sim \log\lp \frac{\exp(-t^2/2)}{\exp(-(t-\mu)^2/2)}\rp \text{ where } t\sim \cN(0,1)\\
		&\sim \frac{(t-\mu)^2}{2} - \frac{t^2}{2} \text{ where } t\sim \cN(0,1)\\
		&\sim \frac{\mu^2}{2}-\mu t \text{ where } t\sim \cN(0,1)\\
		&= \cN\lp \frac{\mu^2}{2},\mu^2 \rp.
	\end{align*}
	A similar calculation shows that $X=\cN\lp -\frac{\mu^2}{2},\mu^2\rp$
\end{proof}




\begin{proposition}[Laplace Mechanism]
	The PRVs for the privacy curve $\delta\lp\Lap\lp \mu,1\rp||\Lap\lp 0,1\rp\rp$ are:
	$$X=|Z|-|Z-\mu| \text{ and } Y= |Z-\mu|-|Z|$$ where $Z\sim \Lap(0,1).$
\end{proposition}
\begin{proof}
	Let $P=\Lap(\mu,1)$ and $Q=\Lap(0,1)$.
	By Theorem~\ref{thm:PRV_LLRV},
	\begin{align*}
		Y &\sim \log\lp \frac{Q(t)}{P(t)}\rp \text{ where } t\sim Q\\
		&\sim \log\lp \frac{\exp(-|t|)}{\exp(-|t-\mu|}\rp \text{ where } t\sim \Lap(0,1)\\
		&\sim |t-\mu|-|t| \text{ where } t\sim \Lap(0,1)\\
		&= |Z-\mu|-|Z| \text{ where } Z\sim \Lap(0,1).
	\end{align*}
	A similar calculation shows that $X=|Z|-|Z-\mu|$ where $Z\sim\Lap(0,1).$
\end{proof}

\begin{proposition}[$(\eps,\delta)$-DP]
	The PRVs for the privacy curve of a $(\eps,\delta)$-DP algorithm are
	$$X=\begin{cases}
	-\infty & \text{w.p. } \delta\\
	-\eps & \text{w.p. }  \frac{(1-\delta)e^\eps}{e^\eps+1}\\
	\eps & \text{w.p. } \frac{1-\delta}{e^\eps+1},
	\end{cases}$$
	$$Y=\begin{cases}
	-\eps & \text{w.p. } \frac{1-\delta}{e^\eps+1}\\
	\eps & \text{w.p. }  \frac{(1-\delta)e^\eps}{e^\eps+1}\\
	\infty & \text{w.p. } \delta.
	\end{cases}$$
\end{proposition}
\begin{proof}
	It is easy to verify that $Y(t)=e^t X(t)$ for all $t\in \R$. We can also verify that $$\delta(\eps)=\Pr[Y> \eps] - e^\eps \Pr[X > \eps]=\delta.$$ Morever $X=-Y$, therefore the privacy curve $\delta(X||Y)$ is symmetric by Proposition~\ref{prop:finv_prv}, i.e., $\delta(X||Y)=\delta(Y||X)$. These conditions together imply that $X,Y$ are PRVs for the $(\eps,\delta)$-DP curve.
\end{proof}



Note that in the all the above examples, we have $X=-Y$ as the privacy curves are symmetric.



\subsection{Subsampling}
In this section, we calculate the PRVs for a subsampled mechanism given the PRVs for the original mechanism. Given two random variables $P,Q$ and a sampling probability $p\in [0,1]$,  $p\cdot P+(1-p)\cdot Q$ denotes the mixture where we sample $P$ w.p. $p$ and $Q$ w.p. $1-p.$
\begin{proposition}
	\label{prop:mixture_f_p_privacy_rv}
	Let $(X,Y)$ be the PRVs for a privacy curve $\delta(P||Q)$. Let $(X_p,Y_p)$ be the PRVs for $\delta_p=\delta(P||\ p\cdot P+(1-p)\cdot Q)$. Then
	\begin{align*}
		X_p&=\log(1+p(e^X-1)),\\
		Y_p&=\begin{cases}
			\log(1+p(e^Y-1)) \text{ w.p. } p\\
			\log(1+p(e^X-1)) \text{ w.p. } 1-p.
		\end{cases}
	\end{align*}
	The CDFs of $X_p$ and $Y_p$ are given by:
	\begin{align*}
		\CDF_{X_p}(t) &= \begin{cases}
			\CDF_X\lp\log\lp \frac{e^t-(1-p)}{p}\rp \rp &\text{ if } t\ge \log(1-p)\\
			0 &\text{ if } t< \log(1-p)
		\end{cases}\\
		\CDF_{Y_p}(t) &= \begin{cases}
			p\cdot \CDF_Y\lp \log\lp \frac{e^t-(1-p)}{p}\rp \rp + (1-p)\cdot \CDF_X\lp \log\lp \frac{e^t-(1-p)}{p}\rp \rp &\text{ if } t\ge \log(1-p)\\
			0 & \text{ if } t <\log(1-p).
		\end{cases}
	\end{align*}

\end{proposition}


\begin{proof}
By Theorem~\ref{thm:PRV_LLRV}, 
\begin{align*}
	X_p &= \log\lp \frac{pY(t)+(1-p)X(t)}{X(t)}\rp \text{ where } t\sim X\\
	& = \log \lp p e^t + 1-p\rp \text{ where } t \sim X\\
	& = \log \lp p e^X + 1-p\rp.
\end{align*}
Similarly,
\begin{align*}
	Y_p &= \log\lp \frac{pY(t)+(1-p)X(t)}{X(t)}\rp \text{ where } t\sim pY+(1-p)X\\
	& = \log \lp p e^t + 1-p\rp \text{ where } t \sim pY+(1-p)X\\
	& = \begin{cases}
			\log(1+p(e^Y-1)) \text{ w.p. } p\\
			\log(1+p(e^X-1)) \text{ w.p. } 1-p.
		\end{cases} 
\end{align*}
The CDF of $X_p$ is given by:
\begin{align*}
	\Pr[X_p \le t] &= \Pr\lb \log \lp p e^X + 1-p\rp \le t\rb\\
	&= \Pr\lb X \le \log\lp \frac{e^t-(1-p)}{p}\rp \rb
\end{align*}
The CDF of $Y_p$ is given by:
\begin{align*}
	\Pr[Y_p \le t] &= p \Pr\lb \log \lp p e^Y + 1-p\rp \le t\rb + (1-p) \Pr\lb \log \lp p e^X + 1-p\rp \le t\rb\\
	&= p\Pr\lb Y \le \log\lp \frac{e^t-(1-p)}{p}\rp \rb + (1-p)\Pr\lb X \le \log\lp \frac{e^t-(1-p)}{p}\rp \rb.
\end{align*}
\end{proof}
\input{appendix}



\end{document}
